\section*{الفصل \ref{ch:b1:elementary-steps} --- آليات التفاعل: الخطوات الأولية والتقريبات}

\begin{solution}{exo:b1:elementary-steps:1}
(أ) أحادية الجزيئية: يمكن أن تكون أولية. (ب) ثنائية الجزيئية: يمكن أن تكون أولية.
(ج) ثلاثة جزيئات وأربع روابط يعاد بناؤها دفعة واحدة: ليست أولية.
(د) ثلاثية الجزيئية، و\ce{N2} يحمل الطاقة المحرَّرة بعيدًا: \omterm{def:b1:elementary-steps:elementary-step}{خطوة أولية}
نادرة لكنها حقيقية.
\end{solution}

\begin{solution}{exo:b1:elementary-steps:2}
$v = k[\mathrm{A}][\mathrm{B}]$؛ $v = k[\mathrm{A}]^2$؛ $v = k[\mathrm{A}]$.
فالمعادلة الإجمالية حصيلة، لا وصف لكيفية التقاء
الجزيئات؛ ويجب أن يقاس قانون سرعتها، أو أن يُستنتج من
آلية.
\end{solution}

\begin{solution}{exo:b1:elementary-steps:3}
الوسيط هو التجويف عند 30؛ والحالتان الانتقاليتان هما
القمتان عند 70 و55. وأعلى \omterm{def:b1:elementary-steps:energy-profile}{حالة انتقالية} هي الأولى: فالخطوة
الأولى هي المحدِّدة للسرعة. وهي ماصة للحرارة (من 0 صعودًا إلى 30).
\end{solution}

\begin{solution}{exo:b1:elementary-steps:4}
بالعامل نفسه $10^4$ (\cref{prop:b1:elementary-steps:catalysis})؛
وثابت التوازن لا يتغير؛ ويُبلغ التوازن أبكر بنحو $10^4$
مرة.
\end{solution}

\begin{solution}{exo:b1:elementary-steps:5}
$t_{\max} = \ln(0.30/0.10)/(0.30 - 0.10) = \qty{5.49}{s}$. ثم
$[\mathrm{B}]_{\max} = a_0\frac{0.10}{0.20}(\eu^{-0.549} - \eu^{-1.648}) =
0.5\,a_0(0.577 - 0.192) = 0.192\,a_0$.
\end{solution}

\begin{solution}{exo:b1:elementary-steps:6}
$\dd[\mathrm{I}]/\dd t = k_1[\mathrm{A}] - (k_{-1} + k_2)[\mathrm{I}] = 0$
يعطي $[\mathrm{I}] = k_1[\mathrm{A}]/(k_{-1} + k_2)$ و$v = k_2[\mathrm{I}]
= \frac{k_1k_2}{k_{-1}+k_2}[\mathrm{A}]$. وإذا كان $k_2 \gg k_{-1}$، فإن $v =
k_1[\mathrm{A}]$: الخطوة الأولى هي المحدِّدة للسرعة. وإذا كان $k_2 \ll k_{-1}$،
فإن $v = (k_1/k_{-1})k_2[\mathrm{A}]$: توازن مسبق تليه خطوة
بطيئة.
\end{solution}

\begin{solution}{exo:b1:elementary-steps:7}
المجموع: \ce{2H2O2 -> 2H2O + O2}. \omterm{def:b1:elementary-steps:catalyst}{الحفاز}: \ce{I-} (يُستهلك ثم
يتجدد)؛ والوسيط: \ce{IO-}. والخطوة الأولى البطيئة تحدد
السرعة: $v = k_1[\ce{H2O2}][\ce{I-}]$.
\end{solution}

\begin{solution}{exo:b1:elementary-steps:8}
الخطوة شديدة الصعود: فوفق مسلّمة هاموند تشبه حالتها الانتقالية
الكربوكاتيون. والمتبادِل الذي يخفض طاقة
الكربوكاتيون يخفض طاقة \omterm{def:b1:elementary-steps:energy-profile}{الحالة الانتقالية} بقدر يكاد يكون مماثلًا،
فينخفض الحاجز وتصير الخطوة أسرع.
\end{solution}

\begin{solution}{exo:b1:elementary-steps:9}
$k_{\mathrm B}/k_{\mathrm C} = \exp(20000/RT)$: \num{3.0e3} عند \qty{300}{K}،
و$\exp(4.01) = 55$ عند \qty{600}{K}. فالتسخين يقلل الأفضلية الحركية
للناتج B، ويتيح للنظام، بجعله تكوّن B عكوسًا، أن يبلغ
الناتج C الأكثر استقرارًا.
\end{solution}

\begin{solution}{exo:b1:elementary-steps:10}
$k_1[\mathrm{A}][\mathrm{M}] = (k_{-1}[\mathrm{M}] + k_2)[\mathrm{A}^*]$،
فيكون $v = k_2[\mathrm{A}^*] = \dfrac{k_1k_2[\mathrm{A}][\mathrm{M}]}{k_{-1}[\mathrm{M}]
+ k_2}$. وتحت ضغط عالٍ ($k_{-1}[\mathrm{M}] \gg k_2$)، $v =
(k_1k_2/k_{-1})[\mathrm{A}]$: الرتبة 1. وتحت ضغط منخفض، $v =
k_1[\mathrm{A}][\mathrm{M}]$: الرتبة 2، إذ يصير التنشيط بالاصطدام
الخطوة البطيئة.
\end{solution}

\begin{solution}{exo:b1:elementary-steps:11}
البرهان هو برهان المبرهنة. ومن أجل $k_1 = k_2 = k$:
$\frac{\dd}{\dd t}([\mathrm{B}]\eu^{kt}) = ka_0$، فيكون $[\mathrm{B}]\eu^{kt} =
ka_0t$ و$[\mathrm{B}] = a_0kt\,\eu^{-kt}$، الذي يبلغ قيمته العظمى عند $t = 1/k$.
\end{solution}

\begin{solution}{exo:b1:elementary-steps:12}
$\mathrm{A} + \mathrm{B} \rightleftharpoons \mathrm{I} + \mathrm{C}$ (سريعة،
ثابتها $K_1$)، ثم $\mathrm{I} + \mathrm{B} \to \mathrm{P}$ (بطيئة،
$k_2$). المجموع هو التفاعل الإجمالي؛ و$[\mathrm{I}] =
K_1[\mathrm{A}][\mathrm{B}]/[\mathrm{C}]$ و$v = k_2[\mathrm{I}][\mathrm{B}]
= k_2K_1[\mathrm{A}][\mathrm{B}]^2/[\mathrm{C}]$: فالناتج الثانوي، المتحرر
في التوازن المسبق، يبطئ التفاعل.
\end{solution}

\begin{solution}{pb:b1:elementary-steps:1}
\textbf{1.} مضاعفة $[\ce{NO}]_0$ تضرب $v_0$ في 4: الرتبة 2؛ ومضاعفة
$[\ce{O2}]_0$ تضربها في 2: الرتبة 1.
\textbf{2.} $v = k[\ce{NO}]^2[\ce{O2}]$؛ $k = 1.0 \times 10^{-5}/((1.0
\times 10^{-3})^2 \times 1.0 \times 10^{-3}) =
\qty{1.0e4}{L^2.mol^{-2}.s^{-1}}$.
\textbf{3.} الاصطدام المتزامن لثلاثة جزيئات نادر، والخطوة
الواحدة لا يمكن أن تصير أبطأ بالتسخين (السؤال 5).
\textbf{4.} $E_a = R\ln(k_{400}/k_{300})/(1/300 - 1/400) = 8.314 \ln(0.1)
/(8.33 \times 10^{-4}) = \qty{-23}{kJ/mol}$.
\textbf{5.} تجري \omterm{def:b1:elementary-steps:elementary-step}{الخطوة الأولية} عبر اصطدامات تعبر
حاجزًا؛ وكسر الاصطدامات القادرة على ذلك، $\eu^{-E_a/RT}$ مع $E_a \ge
0$، لا يمكن إلا أن يزداد مع درجة الحرارة.
\textbf{6.} $-\dd[\ce{NO}]/\dd t = 2v$، $-\dd[\ce{O2}]/\dd t = v$.
\textbf{7.} \ce{2NO -> N2O2} مع \ce{N2O2 + O2 -> 2NO2} يعطيان
\ce{2NO + O2 -> 2NO2}؛ والوسيط هو \ce{N2O2}.
\textbf{8.} كلتاهما ثنائية الجزيئية (وعكس الأولى أحادي الجزيئية).
\textbf{9.} $v = k_2[\ce{N2O2}][\ce{O2}]$.
\textbf{10.} $[\ce{N2O2}] = K_1[\ce{NO}]^2$ مع $K_1 = k_1/k_{-1}$.
\textbf{11.} $v = k_2K_1[\ce{NO}]^2[\ce{O2}]$: وهو القانون المرصود.
\textbf{12.} $k = k_1k_2/k_{-1}$.
\textbf{13.} $\dd[\ce{N2O2}]/\dd t = k_1[\ce{NO}]^2 - k_{-1}[\ce{N2O2}] -
k_2[\ce{N2O2}][\ce{O2}]$.
\textbf{14.} بمساواته بالصفر: $[\ce{N2O2}] = k_1[\ce{NO}]^2/(k_{-1} +
k_2[\ce{O2}])$.
\textbf{15.} $v = k_2[\ce{N2O2}][\ce{O2}]$ يعطي القانون المذكور.
\textbf{16.} $k_{-1} \gg k_2[\ce{O2}]$: يتفكك المثنوي أكثر بكثير
مما يلتقي ثنائي الأكسجين، فتبقى الخطوة الأولى في \omterm{def:b1:extent-q-and-k:quantitative}{حالة توازن}.
\textbf{17.} إذا كان $k_2[\ce{O2}] \gg k_{-1}$، فإن $v = k_1[\ce{NO}]^2$: الرتبة 2 بالنسبة إلى
\ce{NO} و0 بالنسبة إلى \ce{O2}.
\textbf{18.} الرتبة 1 المقيسة بالنسبة إلى \ce{O2}: الحالة الحدية الأولى.
\textbf{19.} $\ln k = \ln k_1 + \ln k_2 - \ln k_{-1}$.
\textbf{20.} بالاشتقاق بالنسبة إلى $T$ واستعمال $\dd\ln k_i/\dd T
= E_i/RT^2$: $E_a = E_1 + E_2 - E_{-1}$.
\textbf{21.} يسرّع التسخين الخطوة البطيئة قليلًا ($E_2$ صغيرة) لكنه يزيح
التوازن المسبق بقوة رجوعًا نحو \ce{NO} ($E_{-1}$ كبيرة): فتقل
المثنويات، وتهبط السرعة الصافية.
\textbf{22.} يقع المثنوي على $E_1 - E_{-1} = \qty{-38}{kJ/mol}$ تحت
المتفاعلات، \omterm{def:b1:elementary-steps:energy-profile}{والحالة الانتقالية} الثانية على $E_2 = \qty{15}{kJ/mol}$ فوق
المثنوي: أي \qty{23}{kJ/mol} تحت المتفاعلات. وأعلى نقطة في
المسلك هي \omterm{def:b1:elementary-steps:energy-profile}{الحالة الانتقالية} الأولى، على \qty{2}{kJ/mol} فقط فوقها.
\textbf{23.} $\exp\big(\frac{23000}{8.314}(\frac{1}{400} - \frac{1}{300})\big)
= 0.10$: أبطأ بعشر مرات عند \qty{400}{K}.
\textbf{24.} $E_a = 2 - 40 + 15 = \textbf{\qty{-23}{kJ/mol}}$، وهي القيمة
المقيسة في الجزء الأول: فالآلية تفسر \omterm{def:b1:rate-laws:order}{قانون السرعة} وتعلقه الغريب
بدرجة الحرارة معًا.
\end{solution}
